\scanpage{22}

% letterhead
\deleted{M. DUFOUR\\
VARITYPISTE\\
131, RUE LECOURBE\\
PARIS XV\textsuperscript{e}\\
$\blacksquare$\\
C. C. P. Paris 14.143-34\\
R. M. 140.029}

$A$ algèbre \inserted{locale noeth. complète,} f\textsuperscript{t} simple sur $k$ pour sa topologie habituelle, alors $A$ est-il f\textsuperscript{t} simple sur $\deleted{\unsure{A|K}}k$ pour sa topologie discrète ???

\deleted{1\textsuperscript{o}) Cas $p>0$} \unsure{On va|On} examiner le cas où $k$ est le corps premier (plus gén. un corps \uline{parfait}), où \unsure{=|ou} où $K/k$ est séparable, $K/k$ l'ex. résiduelle

\deleted{1\textsuperscript{o}) \uline{Cas $p>0$}} Alors \unsure{on a|on}. $A\simeq K[[t_1,\ldots,t_n]]$, $K/k$ une extension \uline{séparable}. \deleted{$K$ est form\textsuperscript{t} simple sur $k$, il suffirait que $A$ le soit sur $K$.}

1\textsuperscript{o}) \uline{Cas $p>0$} \deleted{Alors $A$ est un anneau} \deleted{Si $K$ est parfait, alors} Supposons que $K/k$ ait une $p$-base \uline{finie}, i.e. \margin{$\Omega^{\unsure{1|\vphantom{1}}}_{K/k}$ est fini sur $K$.} \deleted{$A$, ayant une $p$-base \inserted{$(t_1,\ldots,t_n)$} sur $K$, est f\textsuperscript{t} simple sur $K$. \inserted{Comme $K$ est parfait,} il suffit de connaître une $p$-base de $A$ sur $k$. Il en existe si $K$ est fini sur $K^p$, on prend une $p$-base $(x_1\ldots x_n)$ de $K$, et y adjoint $t_1,\ldots,t_n$.} En la complétant par les $t_i$, on trouve une $p$-base de $A$ sur $k$, et $k(A^p)\simeq k\otimes_{k^p}A^p$. [cf ``Diff.'']. Donc dans ce cas on gagne.
% archive sheet number: 21
